\scanpage{29}
\mspage{14}

La \unsure{donnee|donnée} de $(M_0,I)$ équivaut, par le \S~1, à la donnée d'un ens. à quatre éléments $I$, on aura $M_0\simeq V_I$, $M_0^*\simeq\mathcal P_{2,2}(I)$.

\deleted{\unsure{Pour un|Par ailleurs}} \inserted{D'autre} part, \deleted{\unsure{il revient au|à un ens de|la donnée d'un}} \inserted{la donnée d'un} ens à deux éléments $C$, et d'un torseur $S$ sous $\mathbb F_2^C$, équivaut à la donnée d'un carré $Q$, $S$ étant l'ens. des sommets de $Q$, $C$ l'ens de ses codiagonales.

Donc le \deleted{catégorie} \inserted{groupoïde} du deuxième mb de (18) est équivalent aussi à celui des triples $(I,Q;k)$, où $I$ est un ens à quatre éléments, $Q$ un carré, et $k$ un isom.

\[
 k:\underset{\substack{\text{ens des}\\\text{orientations}\\\text{de }I}}{\underline{\underline\omega(I)}}\simeq
 \underset{\substack{\text{ens. des}\\\text{codiagonales}\\\text{de }Q}}{\underline{C(Q)}}. \tag*{(20)}
\]

On trouve donc

\uline{Corollaire 2} On a une équivalence can. de groupoïdes

\[
\underset{\mathbb Z/4\mathbb Z}{\underline\Lambda}\text{-modules }\underset{\text{de rg }2}{\text{libres}}
\xrightarrow{\approx}\text{triplets }(I,Q,k),\quad\text{où} \tag*{(21)}
\]

\[
\left\{\begin{aligned}
&I\quad\text{ens. à }4\text{ élément}\\
&Q\quad\text{carré combinatoire}\\
&k:\underline\omega(I)\simeq C(Q)
\end{aligned}\right.
\]

\uline{Corollaire 3} Si $M$ est associé à $I,Q$, on a un diagramme cartésien canonique

\diagram{\begin{tikzcd}[ampersand replacement=\&,column sep=small]
\&\mathrm{Gl}_{\Lambda}(M)\arrow[dl]\arrow[dr]\&\\
{(22)}\quad\mathfrak S_I\arrow[dr,"\mathrm{sg}"']\&\&\mathcal D_Q\arrow[dl,swap,"{\chi_C\;(\text{opération sur l'ens des codiagonales})}"]\\
\&\{\pm1\}\&
\end{tikzcd}}
% archive sheet number: 28
